\documentclass[11pt]{article}

\usepackage[margin=1.15in]{geometry}
\usepackage{amsmath,amssymb,amsthm}
\usepackage{booktabs}
\usepackage{microtype}
\usepackage{xcolor}
\usepackage[colorlinks=true,linkcolor=blue!50!black,citecolor=blue!50!black,urlcolor=blue!50!black]{hyperref}

\newtheorem{theorem}{Theorem}[section]
\newtheorem{lemma}[theorem]{Lemma}
\newtheorem{corollary}[theorem]{Corollary}
\newtheorem{proposition}[theorem]{Proposition}
\newtheorem{definition}[theorem]{Definition}
\newtheorem{remark}[theorem]{Remark}
\newtheorem{conjecture}[theorem]{Conjecture}
\newtheorem{question}[theorem]{Question}

\newcommand{\N}{\mathbb{N}}
\newcommand{\Z}{\mathbb{Z}}
\newcommand{\Ztwo}{\Z_2}
\newcommand{\Zthree}{\Z_3}
\newcommand{\code}[1]{\texttt{#1}}
\newcommand{\jodd}[2]{j_{#1}(#2)}

\title{The 3-adic Shadow of a Collatz Orbit:\\
       Divergence as Sustained Compression}
\author{M.~Sharpe\thanks{\texttt{msharpe@irendi.com}. The mathematics,
  formalization, and text of this paper were produced in collaboration
  with Claude (Anthropic). All formal results are verified by the
  Lean~4 kernel; see Section~\ref{sec:formal} for the audit.}}
\date{June 12, 2026}

\begin{document}
\maketitle

\begin{abstract}
Let $T$ be the Terras form of the Collatz map and let $j = \jodd{T}{n}$
count odd steps among the first $T$. The exact orbit identity
$2^T\,T^T(n) = 3^{j} n + d$ has correction term $d$ depending only on
the parity word of $n$; reducing modulo $3^j$ eliminates the $n$-term
entirely, so \emph{the $3$-adic residue of an orbit is a function of
its parity word alone}. We develop the consequences, all machine-verified
in Lean~4: the correction obeys the cocycle law
$d(w_1 w_2) = 3^{j_2} d(w_1) + 2^{|w_1|} d(w_2)$, exhibiting the
integer dynamics as the $\Ztwo \times \Zthree$ skew product in
elementary form; orbits whose values remain $\ge N$ survive without
descent only at odd density $\ge \log 2 / \log\frac{3N+1}{N}$, so
divergent orbits are pinned to the critical density
$\log 2/\log 3$, at which the $3$-adic digit production rate
$\delta \log_2 3$ equals exactly $1$ bit per step; and the residues
attainable by surviving windows of length $T$ number at most
$3^T / 2^{\lfloor 17(T-1)/27\rfloor} \approx 2^{0.955T}$ inside fibers
of more than $2^{T-1}$ classes --- the precise $3$-adic mirror, with
identical constants, of Terras' theorem in finite form. Together these
give a quantitative reduction: a divergent Collatz orbit is exactly an
integer that compresses its own $3$-adic digit stream by
$\approx 4.5\%$, at every scale, forever. Experiments sharpen the
picture: the true shadow entropy lies below the binomial bound; the
full-window shadow nearly determines the half-window shadow (a strong
cross-scale rigidity predicted by the cocycle); and Sturmian parity
words --- forced on slowly divergent orbits --- occupy an even thinner
shadow set. We state the finite and ergodic forms of the remaining
problem and locate it within the $\times 2$/$\times 3$ rigidity
landscape. The conjecture remains open.
\end{abstract}

\section{Introduction}

The Collatz conjecture, in Terras form, concerns
$T(n) = n/2$ ($n$ even), $T(n) = (3n+1)/2$ ($n$ odd), and asserts that
every positive integer reaches $1$. Its two failure modes are
nontrivial cycles --- excluded for all lengths up to $183$ halvings,
conditionally on the verified range, in the companion paper
\cite{companion} --- and \emph{divergent orbits}, which are the subject
here.

The companion paper establishes a corridor for divergence: an orbit
segment avoiding descent must run odd-step density above the critical
value $\log 2/\log 3 = 0.6309\ldots$ while within the $\log_2 n$-bit
information budget of its start; densities just above critical are
achieved by genuine integers at every scale; and no certificate
carrying a bounded amount of information about $n$ can decide who
survives. The present paper asks what survival costs on the
\emph{other} side of the arithmetic --- in base $3$ --- and finds that
the cost is precise, quantitative, and of a recognizable kind: a
divergent orbit must \emph{compress its own $3$-adic digit stream},
forever, by a fixed positive margin.

The mechanism is an exact identity. Writing $\jodd{T}{n}$ for the
number of odd steps among the first $T$ and abbreviating
$j = \jodd{T}{n}$,
\begin{equation}\label{eq:exact}
2^T \, T^T(n) \;=\; 3^{j}\, n + d_T(n),
\end{equation}
where the correction $d_T(n)$ --- an explicit sum of products of
powers of $2$ and $3$ along the orbit --- depends only on the
\emph{parity word} $w$ of $n$ (the odd/even pattern of its first $T$
steps), not on $n$ itself. Reducing \eqref{eq:exact} modulo $3^{j}$
annihilates the term $3^j n$:
\[
T^T(n) \;\equiv\; 2^{-T} d(w) \pmod{3^{j}}.
\]
The $3$-adic residue of the endpoint is a function of the word alone.
We call it the \emph{shadow} of the window.

Counting shadows against words yields the reduction. By the corridor,
a surviving window forces a supercritical word, and supercritical
words of length $T$ number at most
$\sum_{j \ge 0.631T}\binom{T}{j} \approx 2^{H(0.6309)T} =
2^{0.950T}$. But the shadow lives modulo $3^{j}$, and --- by an
identity that strikes us as the heart of the matter ---
\[
\delta \cdot \log_2 3 = 1 \quad\text{exactly at}\quad
\delta = \frac{\log 2}{\log 3}:
\]
at the critical density the orbit \emph{mints exactly one bit of
$3$-adic digit information per step}, while its word supplies at most
$0.95$ bits per step. A divergent orbit must therefore emit $3$-adic
digits carrying at least $5\%$ less information than their length ---
sustained compression --- and this is forced at every scale
simultaneously.

\subsection*{Contributions}

\begin{enumerate}
\item \textbf{The shadow theorems} (Section~\ref{sec:shadow}),
machine-verified: the exact form \eqref{eq:exact} with the recursion
for $d$; word-determinacy of $d$; the shadow congruence; and the
cocycle law
\[
d(w_1 w_2) = 3^{j_2}\, d(w_1) + 2^{|w_1|}\, d(w_2),
\]
which is the $\Ztwo \times \Zthree$ skew product written in integers
and governs how shadows couple across scales.
\item \textbf{The density floor} (Section~\ref{sec:floor}),
machine-verified and to our knowledge not in the formal literature:
if all values of an orbit stay $\ge N$ and the window does not descend,
then $2^T N^{j} \le (3N+1)^{j}$ --- odd density at least
$\log 2 / \log\frac{3N+1}{N}$, tending to critical as $N \to \infty$.
This is the finite form of ``divergent orbits have liminf parity
density $\ge \log 2/\log 3$,'' a statement deliberately \emph{not}
delivered by the worst-case master inequalities of \cite{companion}.
\item \textbf{The confinement count} (Section~\ref{sec:count}),
machine-verified: the shadows of all supercritical windows of length
$T$ form an explicit set of at most $3^T/2^{\lfloor 17(T-1)/27\rfloor}$
elements, while every member's fiber holds more than $2^{T-1}$
residues. The constants $17/27$ and $3^T$ are \emph{identical} to those
in the $2$-adic Terras theorem of \cite{companion}: the two mirrors are
one phenomenon.
\item \textbf{The reduction} (Section~\ref{sec:reduction}): finite and
ergodic statements of ``divergence $=$ sustained compression,''
locating the remaining problem in the joint $\times 2$/$\times 3$
rigidity landscape, with an explicit numeric gap.
\item \textbf{Experiments} (Section~\ref{sec:experiments}): the true
shadow entropy undercuts the binomial bound (the word-to-shadow map has
persistent collisions); the full-window shadow nearly determines the
half-window shadow ($1.09$ prefixes per full value at $2T = 18$) ---
cross-scale rigidity, as the cocycle predicts; and Sturmian words at
near-critical slope occupy an even thinner shadow set ($167$ distinct
shadows from $400$ words, against $400$ of $400$ for unstructured words
of matched density).
\end{enumerate}

Everything in Sections \ref{sec:shadow}--\ref{sec:count} is formalized
in Lean~4 over \textsf{mathlib} \cite{mathlib} with no \code{sorry} and
no axioms beyond standard foundations (Section~\ref{sec:formal}).

\section{Preliminaries}

Notation follows \cite{companion}. For $T, n \in \N$ the parity word of
$n$ at length $T$ is $w = (T^t(n) \bmod 2)_{t < T} \in \{0,1\}^T$, and
$\jodd{T}{n} = \sum_t w_t$. Two facts from \cite{companion} are used
throughout: the \emph{parity bijection} (the word of length $T$ is
determined by, and determines, $n \bmod 2^T$), and the \emph{binomial
law} (exactly $\binom{T}{j}$ residues mod $2^T$ have $j$ odd steps).
We also use the dichotomy of \cite{companion}: if $2^T \le n$ and
$n \le T^t(n)$ for all $t \le T$, then $2^T < 2 \cdot 3^{\jodd{T}{n}}$
(``survivors inside the information budget run supercritical'').

\section{The shadow}\label{sec:shadow}

\begin{definition}[\code{dcoef}]
Set $d_0(n) = 0$ and
\[
d_{T+1}(n) \;=\;
\begin{cases}
3^{\jodd{T}{T(n)}} + 2\, d_T(T(n)) & n \text{ odd},\\[2pt]
2\, d_T(T(n)) & n \text{ even}.
\end{cases}
\]
\end{definition}

\begin{theorem}[\code{terras\_exact\_form}]\label{thm:exact}
$2^T\, T^T(n) = 3^{\jodd{T}{n}}\, n + d_T(n)$ for all $T, n$.
\end{theorem}

\begin{proof}
Induction on $T$ via the exact doubling identities
$2\,T(n) = n$ ($n$ even) and $2\,T(n) = 3n+1$ ($n$ odd).
\end{proof}

The master inequalities of \cite{companion} are the two-sided bounds
$3^j - 2^j \le d_T(n) + (3^j - 2^T) \cdot 0 \le \ldots$ --- more
usefully: Theorem~\ref{thm:exact} is the identity those inequalities
sandwich, and it isolates all $n$-independence in $d$:

\begin{theorem}[\code{dcoef\_modEq}]\label{thm:dword}
If $n \equiv m \pmod{2^T}$ then $d_T(n) = d_T(m)$: the correction is a
function of the parity word.
\end{theorem}

\begin{proof}
Induction on $T$; congruence mod $2^{T+1}$ propagates through one
Terras step to congruence mod $2^T$ (doubling identities again), and
equal parities make the two branches of the recursion agree, using
also the word-determinacy of $\jodd{T}{\cdot}$ from the parity
bijection.
\end{proof}

\begin{theorem}[the shadow congruence; \code{shadow\_modEq}]
\label{thm:shadow}
If $n \equiv m \pmod{2^T}$ then
\[
T^T(n) \;\equiv\; T^T(m) \pmod{3^{\jodd{T}{n}}}.
\]
Equivalently: $T^T(n) \bmod 3^{j}$ depends only on $n \bmod 2^T$,
namely $T^T(n) \equiv 2^{-T} d(w) \pmod {3^j}$.
\end{theorem}

\begin{proof}
Two instances of Theorem~\ref{thm:exact} share $j$ and $d$
(Theorem~\ref{thm:dword}); adding $3^j m$ to one and $3^j n$ to the
other gives
$2^T\,T^T(n) + 3^j m = 2^T\, T^T(m) + 3^j n$, whence
$2^T T^T(n) \equiv 2^T T^T(m) \pmod{3^j}$, and $2^T$ is invertible
modulo $3^j$.
\end{proof}

\begin{theorem}[the cocycle law; \code{dcoef\_add}]\label{thm:cocycle}
For all $S, T, n$, writing $n' = T^S(n)$ and $j_2 = \jodd{T}{n'}$,
\[
d_{S+T}(n) \;=\; 3^{j_2}\, d_S(n) + 2^{S}\, d_T(n').
\]
\end{theorem}

\begin{proof}
Expand $2^{S+T}\,T^{S+T}(n)$ through the two windows by
Theorem~\ref{thm:exact} and cancel, using additivity of $\jodd{}{}$
over concatenated windows.
\end{proof}

\begin{remark}
Theorem~\ref{thm:cocycle} is the $\Ztwo \times \Zthree$ skew product
in integer clothing: the second window's $3$-adic weight $3^{j_2}$
scales the first correction (a $3$-adic shift), the first window's
$2$-adic length $2^S$ scales the second (a $2$-adic shift). It also
explains why shadows at different scales are strongly coupled
(Section~\ref{sec:experiments}): modulo $3^{j_1 + j_2}$, the
full-window shadow contains the first-window correction with
multiplier $3^{j_2}$.
\end{remark}

\section{The density floor}\label{sec:floor}

The dichotomy of \cite{companion} pins survivors to supercritical
density only while inside the $\log_2 n$-bit budget; beyond the
budget, the worst-case inequalities go silent --- necessarily so, by
the no-go theorem. For \emph{divergent} orbits a different,
average-case mechanism takes over: once all values are large, the odd
step multiplies by at most $\frac{3N+1}{2N} \approx \frac32$, and the
accounting closes again.

\begin{theorem}[\code{orbit\_upper\_refined}, \code{density\_floor}]
\label{thm:floor}
Let $N \ge 1$ and suppose $T^t(n) \ge N$ for all $t < T$. Then with
$j = \jodd{T}{n}$,
\[
T^T(n) \cdot 2^T N^{j} \;\le\; n\,(3N+1)^{j};
\]
in particular, if also $n \le T^T(n)$ (no descent), then
\[
2^T N^{j} \le (3N+1)^{j},
\qquad\text{i.e.}\qquad
\frac{j}{T} \;\ge\; \frac{\log 2}{\log\frac{3N+1}{N}}
\;\xrightarrow[N \to \infty]{}\; \frac{\log 2}{\log 3}.
\]
\end{theorem}

\begin{proof}
Induction; the odd step costs exactly the inequality
$N(3n+1) \le n(3N+1)$, which is $N \le n$.
\end{proof}

\begin{corollary}\label{cor:liminf}
A divergent orbit (one with $T^t(n) \to \infty$, hence never returning
below any fixed $N$ after some time, and in particular never descending
below suitable re-anchored starting points) has liminf parity density
$\ge \log 2/\log 3$. At that density the orbit performs
$\delta \log_2 3 \ge 1$ \emph{bit of $3$-adic digit production per
step}, with equality exactly on the critical line.
\end{corollary}

\section{The confinement count}\label{sec:count}

\begin{definition}[\code{shadowPair}, \code{shadowSet}]
For a window length $T$ let
$\sigma_T(n) = \bigl(\jodd{T}{n},\, T^T(n) \bmod 3^{\jodd{T}{n}}\bigr)$,
and let $\Sigma_T$ be the set of values $\sigma_T(r)$ over the
supercritical residues $r < 2^T$ (those with
$2^T < 2\cdot 3^{\jodd{T}{r}}$).
\end{definition}

By Theorem~\ref{thm:shadow}, $\sigma_T(n) = \sigma_T(n \bmod 2^T)$ for
every $n$.

\begin{theorem}[\code{surviving\_shadow\_mem},
\code{shadow\_compression}]\label{thm:confine}
If $2^T \le n$ and $n \le T^t(n)$ for all $t \le T$, then
$\sigma_T(n) \in \Sigma_T$ and $2 \cdot 3^{\jodd{T}{n}} > 2^T$.
Moreover
\[
2^{\lfloor 17(T-1)/27 \rfloor} \cdot |\Sigma_T| \;\le\; 3^T .
\]
\end{theorem}

Thus the shadow of every survivor lies in a set of at most
$3^T / 2^{\lfloor 17(T-1)/27\rfloor} \approx 2^{0.9554\,T}$ elements,
while its own fiber $\Z/3^{j}$ offers more than $2^{T-1}$ residues:
an occupancy ratio $\approx 2^{-0.045\,T}$. The proof rides the
binomial law, the image bound $|\Sigma_T| \le \#\{\text{supercritical
residues}\}$, and the integer-only tail estimate
$2^m \sum_{j \ge m}\binom{T}{j} \le 3^T$ of \cite{companion} --- the
\emph{same} estimate, with the \emph{same} threshold $17/27$ from
$3^{17} \le 2^{27}$, that bounds the $2$-adic count of non-descending
integers. Terras' theorem and Theorem~\ref{thm:confine} are mirror
images: one counts who can survive, the other counts where survivors
may stand in base $3$.

\section{The reduction}\label{sec:reduction}

\subsection*{Finite form}

Collatz divergence requires an integer $n$ whose orbit, for every
window length $T$ along a sequence of re-anchored starting points,
(i) runs parity density $\ge \log 2/\log 3 - o(1)$
(Theorem~\ref{thm:floor}), hence mints $3$-adic residue information at
$\ge 1 - o(1)$ bits per step; while (ii) the minted residues carry at
most $\log_2 3 - \tfrac{17}{27} = 0.9554\ldots$ bits per step of
information (Theorem~\ref{thm:confine}), since they are functions of a
supercritical word. \emph{Divergence is a sustained $\approx 4.5\%$
compression of the orbit's own $3$-adic digit stream} ($5.0\%$ at the
entropy-optimal constants), enforced consistently across all scales by
the cocycle law.

\subsection*{Ergodic form}

Embed the orbit diagonally in $\Ztwo \times \Zthree$ and pass to
weak-$*$ limits of empirical measures. The limit is invariant for the
skew product in which the base ($\Ztwo$, conjugate to the parity shift
by the bijection of \cite{companion}) drives the fiber $\Zthree$ by
maps that are $3$-adic isometries on even steps and contractions by
$|3|_3 = \tfrac13$ on odd steps. For a divergent orbit the base
entropy is at most $H(\delta) \le H(0.6309) = 0.950$ bits per step
while the fiber contracts at $\delta \log_2 3 \ge 1$ bit per step, so
the limit measure's $\Zthree$-marginal has (upper) dimension at most
$H(\delta)/(\delta \log_2 3) \le 0.95 < 1$.

\begin{question}[the missing rigidity]\label{q:R}
Must the $\Zthree$-marginal of any weak-$*$ limit of empirical
measures of a divergent \emph{integer} orbit --- a limit whose
\emph{base marginal is necessarily singular}, with parity density
pinned to the supercritical value $\log 2/\log 3$ and entropy at most
$0.95 < 1$ --- have full dimension?
\end{question}

An affirmative answer excludes divergent orbits. We stress that the
singular-base clause is not decoration but the entire content, as a
calibration against the $5x+1$ map shows. For $5x+1$ the critical
density is $\log 2/\log 5 = 0.4307$, \emph{below} the generic $0.5$;
orbits are expected to diverge (the orbit of $7$ famously appears to,
at measured density $\approx 0.49$), and the analogous shadow
machinery applies verbatim: such an orbit sustains an indefinite
$\approx 16\%$ ``compression'' of its $5$-adic stream, with a limit
measure whose $5$-adic marginal has dimension
$\le 1/(0.5 \log_2 5) = 0.86 < 1$. Sustained $p$-adic compression of
an integer orbit is therefore \emph{not} impossible in general; the
naive analogue of Question~\ref{q:R} for $5x+1$ is conjecturally
false. The distinguishing feature of $3x+1$ is that its divergence
requires an \emph{atypical} parity density --- a base measure of
deficient entropy --- where $5x+1$ diverges at the generic,
full-entropy density. Any principle that resolves Collatz must
exploit this singularity of the base, and Question~\ref{q:R} is posed
accordingly. The question sits in the joint $\times 2/\times 3$
tradition of Furstenberg \cite{furstenberg} and Rudolph--Johnson
\cite{rudolph}; we note honestly that even far weaker joint-structure
assertions (the binary digits of $3^k$, say) are open
\cite{senge-straus, stewart}. By the no-go theorem of
\cite{companion}, no bounded-information argument resolves it; by
\code{parity\_pattern\_realized}, the compressed shadow sets are
inhabited at every finite scale, so no counting argument alone can
either.

\section{Experiments}\label{sec:experiments}

Three measurements sharpen the picture beyond the theorems
(script \code{analysis/shadow\_experiments.py}; deterministic seeds).

\subsection*{True shadow entropy}

Restricting to words of length $T$ with exactly
$j = \lceil T \log 2/\log 3\rceil$ odd steps:

\begin{center}
\begin{tabular}{rrrrr}
\toprule
$T$ & $j$ & words $\binom{T}{j}$ & distinct shadows &
$\tfrac{1}{T}\log_2(\text{shadows})$ \\
\midrule
12 & 8  & 495     & 325     & 0.695 \\
15 & 10 & 3{,}003   & 1{,}915   & 0.727 \\
18 & 12 & 18{,}564  & 11{,}522  & 0.750 \\
21 & 14 & 116{,}280 & 70{,}489  & 0.767 \\
\bottomrule
\end{tabular}
\end{center}

\noindent
The word-to-shadow map has persistent collisions ($\sim 0.035$ bits
per step below the word-count rate at these lengths): the true
compression exceeds the formalized bound. (Caveat: these are small
$T$; both rates are still rising toward their asymptotes, and we do
not claim the gap survives the limit.)

\subsection*{Cross-scale rigidity}

At $2T = 18$, $j = 12$: the $18{,}564$ supercritical words yield
$11{,}522$ distinct full-window shadows, $226$ distinct half-window
(prefix) shadows --- and only $12{,}522$ distinct \emph{pairs}
(prefix shadow, full shadow), i.e.\ $0.5\%$ of the product bound. The
full shadow determines the prefix shadow up to $1.09$ candidates on
average. As the cocycle law predicts, shadows at nested scales are
nearly hierarchical; a divergent orbit must satisfy them all
\emph{consistently}, a constraint that counting at a single scale does
not see and which we believe is unexplored.

\subsection*{Structured words cluster}

Mechanical (Sturmian) words at slopes just above critical --- the
itineraries forced on slowly divergent orbits by
Theorem~\ref{thm:floor} together with the dichotomy --- produce $167$
distinct shadows from $400$ samples at $T = 40$, against $400$ of
$400$ for unstructured words of matched density. Structured
itineraries inhabit a far thinner shadow set; arithmetic exclusion of
structured divergence should be attacked through the shadow.

\section{The formalization}\label{sec:formal}

Theorems \ref{thm:exact}--\ref{thm:confine} are formalized in
\code{Shadow.lean} ($\sim$430 lines) over the companion development
\cite{companion} (parity bijection, dichotomy, binomial law, tail
estimate). The audit: a clean build with no \code{sorry};
\code{\#print axioms} reports only Lean's standard foundations for all
seven named theorems, and only \code{propext} with \code{Quot.sound}
--- not even choice --- for Theorems \ref{thm:exact} and
\ref{thm:cocycle}. The numerical claims of
Section~\ref{sec:experiments} are independent Python computations; the
identity \eqref{eq:exact} and the word-determinacy of the shadow were
additionally cross-checked on $2 \times 10^4$ random orbits.

\section{Outlook}

Three levers, in the order we would pull them. (1) \emph{Cross-scale
consistency}: formalize and then exploit the near-hierarchy of shadows
--- the cocycle law makes ``the orbit's shadows cohere at all scales''
a concrete system of congruences whose joint solvability over a single
integer has never been studied. (2) \emph{Structured words}: combine
the clustering observation with the substitution structure of Sturmian
words to attack ``no slowly divergent orbits'' --- the thinness of
their shadow sets is exactly what a Mahler-style argument wants.
(3) \emph{Question~\ref{q:R}}: the ergodic statement, for specialists
in $\times 2 \times 3$ rigidity, now with explicit constants and a
fixed dimension gap. The conjecture remains open; what this paper
contributes is a precise statement of \emph{what kind} of miracle a
counterexample would be: an integer that out-compresses the entropy of
its own arithmetic, at every scale, forever.

\subsection*{Acknowledgments}
The formalization, the experiments, and this text were produced in
collaboration with Claude (Anthropic). The Lean community's
\textsf{mathlib} made the formalization routine.

\subsection*{Reproducibility}
\code{lake build} checks every theorem; \code{\#print axioms}
reproduces the audit; \code{analysis/shadow\_experiments.py}
regenerates Section~\ref{sec:experiments} in minutes on commodity
hardware.

\begin{thebibliography}{10}

\bibitem{companion} M.~Sharpe, The Collatz conjecture at the critical
line: machine-verified no-go, density, and cycle theorems, companion
paper and repository, 2026.

\bibitem{furstenberg} H.~Furstenberg, Disjointness in ergodic theory,
minimal sets, and a problem in Diophantine approximation,
\emph{Math. Systems Theory} \textbf{1} (1967), 1--49.

\bibitem{lagarias-book} J.~C.~Lagarias (ed.), \emph{The Ultimate
Challenge: The $3x+1$ Problem}, American Mathematical Society, 2010.

\bibitem{mathlib} The mathlib Community, The Lean mathematical library,
in \emph{CPP 2020}, ACM, 2020, pp.~367--381.

\bibitem{rudolph} D.~J.~Rudolph, $\times 2$ and $\times 3$ invariant
measures and entropy, \emph{Ergodic Theory Dynam. Systems}
\textbf{10} (1990), 395--406.

\bibitem{senge-straus} H.~G.~Senge and E.~G.~Straus, PV-numbers and
sets of multiplicity, \emph{Period. Math. Hungar.} \textbf{3} (1973),
93--100.

\bibitem{stewart} C.~L.~Stewart, On the representation of an integer
in two different bases, \emph{J. reine angew. Math.} \textbf{319}
(1980), 63--72.

\bibitem{tao} T.~Tao, Almost all orbits of the Collatz map attain
almost bounded values, \emph{Forum of Mathematics, Pi} \textbf{10}
(2022), e12.

\bibitem{terras} R.~Terras, A stopping time problem on the positive
integers, \emph{Acta Arithmetica} \textbf{30} (1976), 241--252.

\end{thebibliography}

\end{document}
